\apartado{2.8}{Distribuciones discretas}{distrib}
\bajada{Tres preguntas de planificación, tres procesos distintos, tres fórmulas
distintas. Confundirlas produce números equivocados que parecen razonables — y
los tres procesos aparecen en cualquier servicio que administre tests y agende
entrevistas.}

\definicion{Las tres, y qué las distingue}{%
  \textbf{Binomial}: un número fijo de intentos independientes, con la misma
  probabilidad en cada uno.\par
  \textbf{Poisson}: conteos en un intervalo de tiempo o espacio, con una tasa
  media conocida. No hay «intentos» que contar.\par
  \textbf{Hipergeométrica}: extracción \emph{sin reposición} de un grupo finito
  y chico, donde la probabilidad cambia en cada extracción.}

\minihistoria{El control de calidad industrial, con protocolos}{%
  En una fábrica se toma una muestra de un lote y se cuenta cuántas piezas
  salen fuera de norma. En un servicio de psicología se toma una muestra de los
  protocolos administrados y se cuenta cuántos salen inválidos: mal
  codificados, incompletos, con patrón de respuesta inconsistente.\par
  Es el mismo problema y la misma fórmula. Cambia el sustantivo.}

\ejercicio{0}{¿Qué es un \emph{ensayo de Bernoulli}? Da dos ejemplos de un
  servicio de psicología. Después explica por qué llamar «éxito» a uno de los
  dos resultados puede resultar engañoso.}{}
\renglones[4]
\respuesta{Un experimento con dos resultados posibles. Ej.: un protocolo es
válido o no; un paciente completa el tratamiento o no.}
\solucion{%
  Un experimento con \textbf{exactamente dos resultados posibles}, que por
  convención se llaman éxito y fracaso, con probabilidades $p$ y $1-p$.
  \par
  De un servicio: un protocolo administrado es válido o inválido; un paciente
  completa el tratamiento o lo abandona; una ficha da positivo o negativo.
  \par
  Lo engañoso es que «éxito» es sólo una etiqueta para el resultado que se
  cuenta. Si estás contando protocolos inválidos, el «éxito» es encontrar uno
  inválido — que obviamente no es una buena noticia. Conviene decirlo en voz
  alta antes de que alguien lo lea en un informe.}

\ejercicio{1}{Un servicio administra 25 protocolos por semana. Por experiencia
  sabe que alrededor del \textbf{4\,\%} sale inválido.
  \textbf{Antes de calcular, escribe cuántos protocolos inválidos crees que
  aparecen en una semana típica:} \underline{\hspace{1.8cm}}
  \par\vspace{0.2em}
  Y arriesga también qué probabilidad hay de que no aparezca
  \emph{ninguno}: \underline{\hspace{1.8cm}}}{Predice primero}
\renglones[4]
\respuesta{El promedio es 1. La probabilidad de ninguno es 36\,\%, bastante más
alta de lo que casi todos estiman.}
\solucion{%
  El promedio es $E(X)=25\times 0{,}04$, o sea \textbf{1} protocolo inválido
  por semana.
  \par
  La probabilidad de que no aparezca ninguno es \textbf{36\,\%}, y se calcula
  en el ejercicio 5. Casi todo el mundo la estima mucho más baja —alrededor del
  10\,\%— porque con un promedio de 1 cuesta aceptar que «cero» sea tan
  frecuente.
  \par
  Esa intuición equivocada es cara en la práctica: una semana sin protocolos
  inválidos no significa que el procedimiento haya mejorado.}
\enclase{Anota las estimaciones. El contraste en el ejercicio 5 es el momento
más útil del apartado.}

\ejercicio{1}{Para cada situación, decide si los intentos son
  \emph{independientes} y si la probabilidad $p$ se mantiene constante.
  \begin{enumerate}[label=(\alph*),topsep=2pt,itemsep=1pt]
    \item Administrar 25 protocolos a personas distintas.
    \item Revisar 3 legajos de una caja de 20, sin devolverlos.
    \item Tamizar a 20 estudiantes elegidos al azar de toda la universidad.
  \end{enumerate}}{}
\renglones[4]
\respuesta{(a) sí y sí; (b) no y no; (c) sí y sí, aproximadamente.}
\solucion{%
  (a) \textbf{Sí}: cada administración es independiente de las otras y $p$ no
  cambia.
  \par
  (b) \textbf{No}: al sacar un legajo sin devolverlo cambia lo que queda, así
  que $p$ cambia en cada extracción. Éste es el caso hipergeométrico.
  \par
  (c) \textbf{Sí, aproximadamente}: técnicamente se muestrea sin reposición,
  pero la población es tan grande frente a 20 que $p$ prácticamente no se mueve.
  \par
  \textbf{La regla práctica}: si la muestra es menos del 10\,\% de la
  población, la binomial es buena aproximación aunque no haya reposición. En
  (b) es 3 de 20, o sea el 15\,\%: no alcanza.}

\ejercicio{2}{Dos preguntas sobre el mismo servicio que exigen fórmulas
  distintas. Identifica cuál corresponde a cada una y explica por qué, sin
  calcular todavía.
  \begin{enumerate}[label=(\alph*),topsep=2pt,itemsep=1pt]
    \item De los 25 protocolos de esta semana, ¿cuántos saldrán inválidos?
    \item ¿Cuántas consultas de urgencia llegarán la semana que viene?
  \end{enumerate}}{Dos casos casi iguales}
\renglones[5]
\respuesta{(a) binomial; (b) Poisson. En (b) no hay un número fijo de intentos.}
\solucion{%
  (a) \textbf{Binomial}: hay un número fijo de intentos (25 protocolos), cada
  uno sale válido o inválido, con la misma probabilidad.
  \par
  (b) \textbf{Poisson}: no hay número fijo de intentos. No se puede decir
  «cuántas consultas de urgencia \emph{no} ocurrieron» — sólo hay una tasa por
  unidad de tiempo.
  \par
  \textbf{La pregunta que las separa}: ¿puedo contar los fracasos? Con los
  protocolos sí: los 24 que salieron bien. Con las urgencias no: las urgencias
  que no pasaron no existen como casos. Si no puedes contar los fracasos, no es
  binomial.}

\ejercicio{2}{Un coordinador calcula así la probabilidad de encontrar
  exactamente 2 protocolos inválidos entre 25, con $p=0{,}04$:
  \begin{quote}\itshape
  $P(X=2) = (0{,}04)^2 \times (0{,}96)^{23} = 0{,}000624$
  \end{quote}
  \textbf{Señala qué le falta} al procedimiento, explica por qué, y da el
  resultado correcto.}{Encuentra el error}
\renglones[6]
\respuesta{Falta el coeficiente $\binom{25}{2}=300$. El correcto es
$\approx 0{,}1877$.}
\solucion{%
  Lo que calculó es la probabilidad de \textbf{una secuencia concreta}: que los
  dos primeros salgan inválidos y los 23 restantes válidos, en ese orden exacto.
  \par
  Pero la pregunta no dice \emph{cuáles} dos. Los dos inválidos pueden ser
  cualquier par de los 25, y hay $\binom{25}{2}=300$ pares posibles:
  \[ P(X=2)=\binom{25}{2}(0{,}04)^{2}(0{,}96)^{23}
     = 300 \times 0{,}000624 \approx 0{,}1877 \]
  \par
  El resultado correcto es 18,8\,\%, unas 300 veces el número que obtuvo.
  \par
  \textbf{Cómo detectarlo}: si una probabilidad de algo perfectamente normal te
  da 0,0006, sospecha. Y fíjate que el coeficiente binomial es justamente la
  combinatoria del apartado 2.4 haciendo su trabajo aquí.}
\enclase{Pregunta qué cuenta el 300. Si pueden decir «los pares de posiciones
donde pueden caer los inválidos», la fórmula dejó de ser un símbolo.}

\ejercicio{3}{Sobre el mismo control de calidad: 25 protocolos por semana,
  $p=0{,}04$ de salir inválido.
  \begin{enumerate}[label=(\alph*),topsep=2pt,itemsep=1pt]
    \item ¿Probabilidad de que \emph{ninguno} salga inválido?
    \item ¿Probabilidad de que salga \emph{al menos uno}?
    \item Compara (a) con tu estimación del ejercicio 2.
  \end{enumerate}}{}
\renglones[6]
\respuesta{(a) $\approx 0{,}3604$. (b) $\approx 0{,}6396$.}
\solucion{%
  (a) \[ P(X=0)=\binom{25}{0}(0{,}04)^{0}(0{,}96)^{25}=(0{,}96)^{25}
     \approx 0{,}3604 \]
  Un 36\,\%: más de una semana de cada tres no tiene ningún protocolo inválido.
  \par
  (b) Por complemento, que es mucho más rápido que sumar desde $X=1$ hasta
  $X=25$:
  \[ P(X\ge 1)=1-P(X=0)\approx 1-0{,}3604=0{,}6396 \]
  \par
  (c) Casi todo el mundo estima el (a) alrededor del 10\,\%. \textbf{La
  consecuencia práctica es real}: si el servicio interpreta una semana limpia
  como señal de que el procedimiento mejoró, se está engañando — pasa una de
  cada tres semanas sin que cambie nada.}

\ejercicio{3}{Dos preguntas más del mismo servicio:
  \begin{enumerate}[label=(\alph*),topsep=2pt,itemsep=1pt]
    \item Recibe en promedio 5 solicitudes de atención por semana.
      ¿Probabilidad de recibir exactamente 8 en una semana?
    \item De los 43 estudiantes que dieron positivo, 9 tienen el expediente
      incompleto. Se auditan 6 al azar, sin reposición. ¿Probabilidad de que
      exactamente 2 estén incompletos?
  \end{enumerate}}{}
\renglones[7]
\respuesta{(a) $\approx 0{,}0653$. (b) $\approx 0{,}2739$.}
\solucion{%
  (a) \textbf{Poisson}, con $\lambda = 5$:
  \[ P(X=8)=\frac{e^{-5}\,5^{8}}{8!}
     = \frac{0{,}006738\times 390\,625}{40\,320}\approx 0{,}0653 \]
  \par
  (b) \textbf{Hipergeométrica}: se audita sin reposición de un grupo finito de
  43. Con $N=43$, $K=9$, $n=6$, $k=2$:
  \[ P(X=2)=\frac{\binom{9}{2}\binom{34}{4}}{\binom{43}{6}}
     = \frac{36\times 46\,376}{6\,096\,454}\approx 0{,}2739 \]
  \par
  Aquí \emph{no} se puede usar binomial: 6 de 43 es el 14\,\%, muy por encima
  del 10\,\%, así que la probabilidad cambia de manera apreciable en cada
  extracción. Se comprueba en el ejercicio 8.}

\ejercicio{4}{Una auditora resuelve el punto (b) anterior con la binomial,
  usando $n=6$ y $p=9/43=0{,}209$, y obtiene $0{,}2569$.
  \begin{enumerate}[label=(\alph*),topsep=2pt,itemsep=1pt]
    \item ¿Por cuánto se equivocó, en términos relativos?
    \item ¿Por qué la binomial no corresponde aquí?
    \item ¿En qué caso su atajo \emph{sí} habría sido aceptable?
  \end{enumerate}}{}
\renglones[7]
\respuesta{(a) $0{,}2569$ contra $0{,}2739$: un 6\,\% menos. (b) Hay extracción
sin reposición de un grupo chico. (c) Si el grupo fuera mucho más grande, por
ejemplo auditar 6 de 500.}
\solucion{%
  (a) Obtuvo $0{,}2569$ contra el valor correcto $0{,}2739$: se equivocó en
  torno al 6\,\% relativo.
  \par
  (b) La binomial supone que $p$ es constante. Aquí no lo es: si el primer
  legajo sale incompleto, quedan 8 incompletos entre 42, y $p$ baja de
  $0{,}209$ a $0{,}190$. Con 6 de 43 esa variación ya pesa.
  \par
  (c) Si el grupo fuera mucho más grande —auditar 6 de 500, por ejemplo— la
  probabilidad casi no se movería y la binomial sería una aproximación
  perfectamente aceptable.
  \par
  \textbf{Lo que hay que retener no es el 6\,\%.} Es que el error es
  \emph{chico}: no salta a la vista, no rompe nada, y por eso sobrevive hasta
  el informe final. Los errores grandes se detectan solos; éstos no.
  \par
  \textbf{El árbol de decisión}: ¿hay un número fijo de intentos? Si no,
  Poisson. Si sí, ¿extraes sin reposición de un grupo chico? Si sí,
  hipergeométrica. Si no, binomial.}
\enclase{Si tienes tiempo, pídeles que calculen cómo cambia $p$ tras la primera
extracción. Ver el paso de $0{,}209$ a $0{,}190$ convence más que la regla del
10\,\%.}

\trampa{usar la binomial para auditorías y muestreos de archivos}%
  {es la fórmula más conocida, y el resultado se parece tanto al correcto que
   nadie lo cuestiona}%
  {en una auditoría se extrae sin reposición de un grupo finito. Si la muestra
   supera el 10\,\% del grupo, corresponde la hipergeométrica}

\cierre{%
  Las tres distribuciones responden preguntas que un servicio se hace todas las
  semanas: cuántos protocolos van a fallar, cuántas urgencias van a llegar,
  cuántos legajos incompletos va a encontrar una auditoría.
  \par
  Las tres trabajan con variables discretas — cosas que se cuentan. Lo que
  falta es la herramienta para las que se miden, y con ella se puede evaluar
  cualquier punto de corte sin volver a contar el archivo.}
